%% notes-tex/w037-isobutanol-scrnaseq/sections/3-selection.tex
%% [[experiments.W037-isobutanol-scrnaseq.strain-selection]]
%% https://github.com/Mjvolk3/torchcell/tree/main/notes-tex/w037-isobutanol-scrnaseq/sections/3-selection.tex
%%
%% SOURCE: experiments/W037-isobutanol-scrnaseq/scripts/strain_tables.py
%% Tables 2 and 3 are generated by that script. The fold changes quoted in the
%% prose are the derivation fields of the Titer records behind Table 2, and the
%% deletion lists are the deletions fields of the same records.

\section{The six to grow\secstatus{todo}}\label{sec:selection}

Six of the eight isobutanol strains ferment 15\% glucose, so one protocol covers
them. They form three pairs, and within a pair the two strains share a background
and differ only by the plasmid they carry.

%% GENERATED FILE -- do not hand-edit.
%% SOURCE: experiments/W037-isobutanol-scrnaseq/scripts/strain_tables.py
\begin{table}[H]\centering
\footnotesize
\caption[The six selected]{The six selected, as three axes of one variable
each. Within an axis the two strains share a background and differ by the
plasmid they carry, so a difference measured between them is attributable to
that plasmid. \emph{pathway} names the compartment the heterologous
isobutanol route was built in. Every strain here ferments 15\% glucose, so all
six run under one protocol. Titer conventions follow
Table~\ref{tab:inventory}.}\label{tab:selected}
\begin{tabular}{lllllll}
\toprule
axis & id & strain & deletions & pathway & plasmid & titer (mg/L) \\
\midrule
\multirow{2}{*}{1} & JC0052 & YZy91 & \textit{bat1}$\Delta$, \textit{bat2}$\Delta$, \textit{ilv6}$\Delta$ & -- & none & not reported \\
 & JC0043 & YZy311 & \textit{bat1}$\Delta$, \textit{bat2}$\Delta$, \textit{ilv6}$\Delta$ & -- & ILV6 WT & $\sim$118\,$^{\dagger}$ \\
\addlinespace
\multirow{2}{*}{2} & JC0044 & YZy313 & \textit{bat1}$\Delta$ & mitochondrial & ILV6 WT & not reported \\
 & JC0045 & YZy314 & \textit{bat1}$\Delta$ & mitochondrial & ILV6 V110E & 1.6$\times$ JC0044 \\
\addlinespace
\multirow{2}{*}{3} & JC0050 & YZy454 & \textit{ilv3}$\Delta$, \textit{tma29}$\Delta$ & cytosolic & Ll\_ilvD WT & $\sim$73\,$^{\dagger}$ \\
 & JC0051 & YZy469 & \textit{ilv3}$\Delta$, \textit{tma29}$\Delta$ & cytosolic & Ll\_ilvD I433V & \textbf{227 $\pm$ 13} \\
\bottomrule
\end{tabular}
\end{table}


\notesec{Axis 1 is a deletion against its add-back}
\strain{JC0052} is \gene{bat1}$\Delta$ \gene{bat2}$\Delta$ \gene{ilv6}$\Delta$ with
the biosensor and no pathway genes. \strain{JC0043} is that same strain carrying
\gene{ILV6} on a CEN plasmid. \gene{ILV6} encodes the regulatory subunit of
acetolactate synthase, and deleting it relieves valine inhibition, which is why
the deletion raises isobutanol in the first place. This pair is the most direct
point of contact with knockout-collection data, because the lesion is a single
gene whose isobutanol phenotype a deletion screen reports on.

\notesec{Axis 2 is one allele of one gene, in a producing background}
\strain{JC0044} and \strain{JC0045} are both YZy363, which is \gene{bat1}$\Delta$
with a five-cassette mitochondrial isobutanol pathway $\delta$-integrated and
\gene{ILV6} intact. They differ in the plasmid allele, wild type against V110E,
and the paper reports V110E as 1.6 times the wild-type titer without giving
either absolute. Because \gene{ILV6} is intact here, both strains overexpress the
gene on top of a native copy, which is a different genetic situation from Axis 1
and not a repeat of it.

\notesec{Axis 3 is a heterologous enzyme against its improved variant}
\strain{JC0050} and \strain{JC0051} are both YZy452, which is \gene{ilv3}$\Delta$
\gene{tma29}$\Delta$ with a cytosolic pathway built from
\gene{Bs\_alsS}, \gene{Ec\_ilvC}\gsup{P2D1-A1} and \gene{Ll\_ilvD}. They differ
in the plasmid allele of \gene{Ll\_ilvD}, wild type against I433V, and the
variant is 3.1 times the wild type at 227\,$\pm$\,13\,mg/L. This axis carries the
largest measured difference of the three and the highest glucose-grown titer in
the set. Its two deletions, \gene{ILV3} and \gene{TMA29}, are both genes a
deletion screen covers, so it is a second point of contact with knockout data
rather than a purely synthetic contrast.

%% GENERATED FILE -- do not hand-edit.
%% SOURCE: experiments/W037-isobutanol-scrnaseq/scripts/strain_tables.py
\begin{table}[H]\centering
\footnotesize
\caption[Genotypes]{Genotypes of the six, as Supplementary Table 1 of the
source paper gives them, with the selection markers and the growth medium each
one needs. \emph{medium} follows from the plasmid: a strain carrying a
URA3 plasmid is held on SC-ura, and JC0052, which carries none, needs uracil
instead. Section~\ref{sec:before} treats that mismatch, which is the one
thing in this list that has to be fixed before the strains are
grown.}\label{tab:genotypes}
\begin{tabular}{@{}l p{0.86\textwidth}@{}}
\toprule
id & strain and genotype \\
\midrule
\textbf{JC0052} & YZy91 \\
 & {\scriptsize \textit{his3}::HIS3-P\psub{LEU1}-yEGFP-PEST-T\psub{ADH1}-P\psub{TPI1}-LEU4\gsup{1-410}-T\psub{PGK1}, \textit{bat1}$\Delta$::hphMX, \textit{bat2}$\Delta$::lox71-kanMX-lox66, \textit{ilv6}$\Delta$::lox71-natMX-lox66} \\
 & {\scriptsize markers: hphMX, kanMX, natMX, HIS3 \quad medium: SC + uracil} \\
\addlinespace
\textbf{JC0043} & YZy311 \\
 & {\scriptsize YZy91, CEN URA3 plasmid (P\psub{TDH3}-ILV6-T\psub{ADH1})} \\
 & {\scriptsize markers: hphMX, kanMX, natMX, HIS3, URA3 \quad medium: SC-ura} \\
\addlinespace
\textbf{JC0044} & YZy313 \\
 & {\scriptsize YZy363 ($\delta$-integrated P\psub{TDH3}-ILV2, P\psub{PGK1}-ILV3, P\psub{TEF1}-CoxIVMLS-Ll\_adhA\gsup{RE1}, P\psub{TDH3}-CoxIVMLS-ARO10, P\psub{TEF1}-ILV5), CEN URA3 plasmid (P\psub{TDH3}-ILV6-T\psub{ADH1})} \\
 & {\scriptsize markers: hphMX, HIS3, ShBle, URA3 \quad medium: SC-ura} \\
\addlinespace
\textbf{JC0045} & YZy314 \\
 & {\scriptsize YZy363, CEN URA3 plasmid (P\psub{TDH3}-ILV6\gsup{V110E}-T\psub{ADH1})} \\
 & {\scriptsize markers: hphMX, HIS3, ShBle, URA3 \quad medium: SC-ura} \\
\addlinespace
\textbf{JC0050} & YZy454 \\
 & {\scriptsize YZy452, CEN URA3 plasmid (P\psub{TDH3}-Ll\_ilvD-T\psub{ADH1})} \\
 & {\scriptsize markers: hphMX, kanMX, HIS3, ShBle, URA3 \quad medium: SC-ura} \\
\addlinespace
\textbf{JC0051} & YZy469 \\
 & {\scriptsize YZy452, CEN URA3 plasmid (P\psub{TDH3}-Ll\_ilvD\gsup{I433V}-T\psub{ADH1})} \\
 & {\scriptsize markers: hphMX, kanMX, HIS3, ShBle, URA3 \quad medium: SC-ura} \\
\bottomrule
\end{tabular}
\end{table}


\notesec{What the three axes buy together}
The six span three pathway states, native flux in Axis 1, mitochondrial in
Axis 2 and cytosolic in Axis 3, three genetic backgrounds, and roughly a
threefold measured range in titer, on one carbon source and one assay format.
Axis 1 carries no engineered route at all: both of its strains make isobutanol
through endogenous valine biosynthesis and the Ehrlich pathway, with
\gene{bat1}$\Delta$ \gene{bat2}$\Delta$ keeping
$\alpha$-ketoisovalerate from being transaminated away. The plasmid that
separates its two strains carries \gene{ILV6}, a regulatory subunit rather than
a pathway enzyme, so the pair varies feedback control and not pathway capacity. Each pair isolates one variable,
so a transcriptional difference within a pair has one candidate cause. Adding
\strain{JC0044} is what makes Axis 2 a pair rather than a single high point, and
without it \strain{JC0045} has no same-background comparator anywhere in the set.
